% ===== Section 5: "You're pretty smart for a small person" (~16 min) =====
\topicopener{giraffe}{Հասակ ու գեներ}{Փոքր մարդու համար բավական խելացի ես}{Հասակ, խելք ու ժառանգականություն}{happy}

% Myth or fact: tall people are smarter
\begin{frame}{Առասպել, թե՞ փաստ \#2}
  \mythorfact{Բարձրահասակներն ավելի խելացի են}{ՓԱՍՏ}{mintdark}{%
    \begin{card}[colback=mint!35!cardwhite]
      Բըրնեթ. «Դա ճիշտ է»։ Բայց \textbf{բացարձակ չէ}. բասկետբոլիստներն ավելի խելացի չեն ժոկեյներից, Հսկա Անդրեն խելացի չէր Այնշտայնից, իսկ Հագրիդը չէր կարող խորամանկել Մարի Կյուրիին։
    \end{card}}
\end{frame}

% How much smarter: correlation 0.2, 3-4 IQ points
\begin{frame}{Որքա՞ն ավելի խելացի}
  \begin{columns}[c]
    \begin{column}{0.5\textwidth}
      \centering
      % Simulated data (220 people, seed 509): height ~ N(172, 9) cm, IQ ~ N(100, 15), correlation 0.2
% (sample r = 0.23). Not real measurements: it only shows how weak a correlation of 0.2 looks.
% The orange line is the r = 0.2 regression line (slope 0.2 * 15/9 IQ points per cm).
\begin{tikzpicture}
  \begin{axis}[
      width=6.6cm, height=5.0cm,
      xmin=145, xmax=200, ymin=55, ymax=145,
      xtick={150,160,170,180,190}, ytick={70,100,130},
      xlabel={հասակ, սմ}, ylabel={IQ},
      label style={font=\small, text=plum}, tick label style={font=\small, text=plum},
      axis lines=left, axis line style={plum},
    ]
    \addplot[only marks, mark=*, mark size=1.3pt, color=amBlue, fill opacity=0.45, draw opacity=0.6] coordinates {(166.4,66.7) (173.3,109.0) (179.9,117.1) (159.1,87.4) (167.9,87.9) (165.1,126.9) (170.5,97.8) (169.5,108.7) (165.6,103.3) (177.3,103.4) (161.4,96.5) (162.0,101.7) (160.8,79.5) (178.1,102.0) (166.9,95.3) (177.2,104.3) (166.5,119.4) (179.6,95.2) (169.8,85.8) (172.9,110.5) (176.8,94.7) (153.8,78.6) (160.6,108.6) (179.1,92.3) (155.6,103.3) (167.3,84.6) (177.3,81.7) (177.4,85.4) (175.8,94.5) (167.5,99.0) (169.8,113.1) (182.6,89.6) (169.1,84.5) (182.3,96.3) (167.4,101.8) (155.1,100.9) (158.6,120.7) (176.7,74.0) (172.3,103.7) (172.5,102.1) (175.8,129.8) (172.5,96.0) (160.9,89.9) (164.4,112.0) (188.8,112.2) (176.8,94.9) (177.3,111.2) (191.6,105.4) (158.7,99.7) (172.7,102.7) (182.1,128.2) (190.8,95.8) (156.6,87.9) (167.7,106.2) (163.8,140.5) (170.5,113.2) (178.5,98.4) (169.8,97.6) (173.5,95.1) (170.9,108.8) (190.1,109.9) (169.6,105.4) (176.7,102.7) (178.8,113.9) (168.6,118.4) (176.7,100.0) (176.3,118.1) (165.8,109.8) (161.4,89.4) (177.3,93.8) (182.6,104.9) (165.3,98.8) (165.3,106.0) (182.4,99.6) (164.9,121.3) (171.0,109.9) (146.2,99.0) (166.8,74.9) (180.8,113.7) (188.0,105.9) (175.3,94.6) (163.8,105.1) (175.3,115.9) (176.5,107.3) (181.8,117.0) (157.1,88.2) (182.4,106.2) (172.3,85.4) (173.6,109.4) (181.3,81.2) (164.2,111.1) (182.7,111.9) (176.0,99.6) (155.0,87.9) (167.6,85.0) (165.7,73.5) (165.2,99.2) (198.0,106.4) (177.4,72.8) (160.2,87.5) (169.6,70.5) (179.6,99.8) (157.4,122.0) (151.0,111.2) (171.7,86.8) (182.0,108.6) (165.8,96.8) (180.6,102.1) (194.3,108.8) (174.5,89.5) (176.2,89.2) (173.6,82.1) (180.7,87.6) (187.4,120.3) (175.9,97.4) (181.8,111.0) (167.9,132.2) (173.6,105.0) (176.3,92.5) (159.1,79.0) (183.2,91.6) (171.3,83.6) (174.4,79.2) (174.3,94.3) (174.7,94.0) (169.4,106.3) (170.8,76.7) (172.1,109.9) (173.4,112.4) (171.8,96.8) (167.8,109.2) (173.0,116.6) (163.1,103.5) (189.2,89.7) (169.4,114.5) (187.0,106.5) (158.2,111.4) (179.1,94.1) (177.6,95.2) (155.5,74.5) (177.3,63.0) (173.3,124.2) (174.8,61.4) (161.6,66.8) (156.4,97.7) (163.9,105.0) (177.4,111.8) (160.2,99.9) (162.8,104.2) (148.6,99.6) (172.0,127.6) (169.2,93.7) (178.5,102.6) (175.4,114.8) (170.3,98.8) (171.6,92.1) (166.9,95.4) (186.0,106.2) (178.9,129.8) (176.5,119.1) (170.3,116.8) (183.7,116.9) (181.1,117.1) (177.6,101.1) (157.5,93.6) (183.9,110.3) (176.5,83.0) (177.2,72.1) (179.8,97.5) (189.6,119.2) (171.1,94.3) (176.9,100.9) (164.5,86.2) (183.2,101.1) (162.9,90.0) (161.8,84.2) (166.9,100.8) (177.8,96.8) (169.9,94.9) (172.4,108.6) (169.9,77.2) (169.6,94.1) (165.7,57.4) (193.7,121.5) (171.7,99.8) (184.0,144.0) (179.7,111.5) (168.3,101.8) (168.3,99.8) (173.5,105.8) (161.8,78.3) (173.8,115.6) (164.1,94.7) (167.9,105.8) (182.4,119.8) (175.0,91.2) (173.8,80.1) (180.1,132.4) (154.8,56.3) (183.3,99.3) (175.7,136.1) (164.7,97.6) (169.0,123.0) (162.0,90.7) (159.9,93.7) (181.5,119.2) (177.1,115.3) (167.4,111.0) (181.8,139.2) (169.3,101.5) (177.4,108.9) (176.4,121.9) (183.7,103.1) (159.2,98.6) (159.4,139.6) (168.7,98.1) (171.8,144.0) (170.6,86.3) (162.9,91.3) (183.8,115.3)};
    \addplot[amOrange, line width=2pt, domain=147:198] {100+0.3333*(x-172)};
  \end{axis}
\end{tikzpicture}
\par
      {\scriptsize\color{plum!70!cream} մոդելավորված տվյալներ, կոռելյացիան՝ 0.2\par}
    \end{column}
    \begin{column}{0.46\textwidth}
      \begin{itemize}
        \item Հասակի ու IQ-ի \textbf{կոռելյացիան} (correlation) մոտ \textbf{0.2} է
        \item<2-> Վերցրու պատահական մեկ բարձրահասակի ու մեկ ցածրահասակի. ով կլինի խելացի, \textbf{անհայտ է}
        \item<3-> Իսկ 10 000-ը 10 000-ի դեմ՝ բարձրահասակների միջին IQ-ն մի փոքր բարձր է, \textbf{3-4 միավորով}
        \item<4-> Արդյունքը կրկնվում է շատ հետազոտություններում
      \end{itemize}
    \end{column}
  \end{columns}
  \source{Taki et al., NeuroImage, 2012}
\end{frame}

% What a correlation of 0.2 means (correction to the book)
\begin{frame}{Ի՞նչ է նշանակում 0.2}
  Կոռելյացիան ցույց է տալիս, թե որքան են երկու բան \textbf{միասին} փոխվում. 0՝ ոչ մի կապ, 1՝ կատարյալ կապ։\par\medskip
  \begin{columns}[c]
    \begin{column}{0.52\textwidth}
      \onslide<2->{%
      \begin{scienceupdate}
        Գրքում՝ «0.2 նշանակում է, որ կապը կա \textbf{5 մարդուց 1-ի} մոտ»։ Այդպես չէ. կոռելյացիան ամբողջ խմբի մասին է, ոչ առանձին մարդկանց։\par\smallskip
        Կոպիտ ասած՝ հասակը «բացատրում» է IQ-ի տարբերությունների ընդամենը \textbf{մոտ 4\%-ը} (0.2 $\times$ 0.2 = 0.04)։
      \end{scienceupdate}}
    \end{column}
    \begin{column}{0.44\textwidth}
      \centering
      \onslide<3->{%
      \begin{tikzpicture}
        \fill[amOrange] (0,0) rectangle (0.24,0.9);
        \fill[amBlue!25!cream] (0.24,0) rectangle (6,0.9);
        \draw[plum, line width=0.9pt] (0,0) rectangle (6,0.9);
        \node[anchor=south west, font=\blackfont\small, text=amOrange!80!black] at (0,0.95) {4\%՝ հասակ};
        \node[font=\blackfont\small, text=plum] at (3.2,0.45) {96\%՝ մնացած ամեն ինչ};
      \end{tikzpicture}\par\medskip
      {\small IQ-ի տարբերությունները մարդկանց միջև\par}}
    \end{column}
  \end{columns}
\end{frame}

% Meme: Roll Safe, correlation is not causation
\begin{frame}{Կատարյալ պլան \#2}
  \centering
  \begin{meme}[10.6cm]{roll_safe.jpg}
    \memecap[text width=8cm, font=\blackfont\normalsize]{0.5}{0.88}{Բարձրահասակներն ավելի խելացի են}
    \memecap[text width=8cm, font=\blackfont\normalsize]{0.5}{0.1}{Ուրեմն քննությանը կգնամ բարձր կրունկներով}
  \end{meme}
\end{frame}

% Heritability: definition with the book's examples
\begin{frame}{Ժառանգելիություն (heritability)}
  \begin{card}[colback=butter!45!cream]
    Ամենահավանական բացատրությունը \textbf{գենետիկ} է։ Ժառանգելիությունը ցույց է տալիս, թե խմբում հատկանիշի \textbf{տարբերությունների} որ մասն է գեներից. 1.0՝ բոլորը, 0.0՝ ոչ մեկը։
  \end{card}
  \vspace{1mm}
  \centering
  % Heritability from 0 to 1 with the book's examples: being on fire (0.0), species (1.0),
% and IQ in children (~0.45) and adults (~0.85) from Bouchard's review. Overlays 2-5.
\begin{tikzpicture}[x=10.5cm]
  \shade[left color=peach, right color=lavender, rounded corners=1.5mm] (0,-0.16) rectangle (1,0.16);
  \draw[plum, line width=0.9pt, rounded corners=1.5mm] (0,-0.16) rectangle (1,0.16);
  \foreach \v in {0,0.5,1}{\node[font=\small, text=plum, anchor=north] at (\v,-0.2) {\v};}
  % ends: the book's two examples
  \node[visible on=<2->, anchor=south, align=center, font=\small, text=plum, text width=3.6cm] at (0.08,0.25)
    {\showon{2-}{\bigemoji[0.8cm]{fire}}\\ Հիմա այրվո՞ւմ ես։ Դա միայն միջավայրից է՝ \textbf{0.0}};
  \node[visible on=<3->, anchor=south, align=center, font=\small, text=plum, text width=3.6cm] at (0.92,0.25)
    {\showon{3-}{\bigemoji[0.8cm]{dna}}\\ Ծնողներդ խոզ են՝ դու էլ խոզ ես. տեսակը՝ \textbf{1.0}};
  % IQ according to Bouchard
  \fill[amBlue, visible on=<4->] (0.45,0) circle (1.6mm);
  \fill[amRed, visible on=<5->] (0.85,0) circle (1.6mm);
  \node[visible on=<4->, anchor=north, font=\blackfont\small, text=amBlue] at (0.45,-0.55) {IQ, երեխաներ $\approx$ 0.45};
  \node[visible on=<5->, anchor=north, font=\blackfont\small, text=amRed] at (0.82,-0.55) {IQ, մեծեր $\approx$ 0.85};
\end{tikzpicture}

\end{frame}

% IQ heritability, children vs adults (+ science update)
\begin{frame}{Ինչու են երեխաների ու մեծերի թվերը տարբեր}
  \begin{columns}[T]
    \begin{column}{0.54\textwidth}
      \begin{itemize}
        \item Թոմաս Բուշարի ակնարկով՝ մեծերի մոտ \textbf{0.85}, երեխաների մոտ \textbf{0.45}
        \item<2-> Գեներն ավելի ուժե՞ղ են մեծերի մոտ։ Ոչ. խոսքը խմբի \textbf{տարբերությունների} մասին է
        \item<3-> Երեխայի ուղեղը դեռ զարգանում է, ու միջավայրը շատ բան է փոխում։ Մեծերը մոտավորապես նույն կրթությունն են ստացել, ու տարբերությունները մնում են գեներին
      \end{itemize}
    \end{column}
    \begin{column}{0.42\textwidth}
      \onslide<4->{%
      \begin{scienceupdate}
        Նոր տվյալներ. 9 տարեկանում մոտ \textbf{0.4}, 17-ում՝ \textbf{0.66}, ավելի ուշ՝ մինչև \textbf{0.8}։ Պատկերը նույնն է, իսկ 0.85-ը վերին սահմանին է։
      \end{scienceupdate}}
    \end{column}
  \end{columns}
  \source{Bouchard, 1997; Haworth et al., Mol. Psychiatry, 2010; Plomin \& Deary, Mol. Psychiatry, 2015}
\end{frame}

% No gene for intelligence (+ GWAS update)
\begin{frame}{Խելքի գեն չկա}
  \begin{columns}[T]
    \begin{column}{0.54\textwidth}
      \begin{itemize}
        \item Ոմանք հույս ունեն, որ կա «խելքի գեն», որը միացնես, ու կխելոքանաս
        \item<2-> Բայց խելքը շատ պրոցեսների գումար է, ու դրանք կառավարում են \textbf{շատ գեներ}
      \end{itemize}
      \vspace{1mm}
      \onslide<3->{%
      \begin{bookquote}
        Հարցնել, թե որ գենն է պատասխանատու խելքի համար, նույնն է, թե հարցնես՝ դաշնամուրի որ ստեղնն է պատասխանատու սիմֆոնիայի համար։ \emoji{musical}
      \end{bookquote}}
    \end{column}
    \begin{column}{0.42\textwidth}
      \centering
      \bigemoji[1.1cm]{dna}\par\medskip
      \onslide<4->{%
      \begin{scienceupdate}
        2018. Մոտ 270 000 մարդու գենոմի վերլուծությունը գտավ \textbf{205} գենոմային գոտի, որոնք կապված են խելքի հետ։ Ամեն մեկի ազդեցությունը \textbf{շատ փոքր} է։
      \end{scienceupdate}}
    \end{column}
  \end{columns}
  \source{Savage et al., Nature Genetics, 2018}
\end{frame}

% Four explanations for the height-IQ link
\begin{frame}{Ինչու՞ են բարձրահասակները խելացի. 4 վարկած}
  \begin{columns}[T]
    \begin{column}{0.47\textwidth}
      \begin{card}[colback=lavender!35!cardwhite]
        \emoji[1.3em]{dna}\ {\blackfont 1. Ընդհանուր գեներ}\par\smallskip
        Մեկ գենը կարող է ազդել և՛ հասակի, և՛ խելքի վրա։ Սա կոչվում է \textbf{պլեյոտրոպիա} (pleiotropy)
      \end{card}
      \onslide<3->{%
      \begin{card}[colback=mint!35!cardwhite]
        \emoji[1.3em]{apple}\ {\blackfont 3. Սնունդ ու առողջություն}\par\smallskip
        Լավ սնունդն ու առողջ մանկությունը մեծացնում են և՛ հասակը, և՛ ուղեղը
      \end{card}}
    \end{column}
    \begin{column}{0.47\textwidth}
      \onslide<2->{%
      \begin{card}[colback=brainpink!25!cardwhite]
        \emoji[1.3em]{heart}\ {\blackfont 2. Սեռական ընտրություն}\par\smallskip
        Բարձրահասակ ու խելացի տղամարդիկ իբր ավելի գրավիչ են ու ավելի շատ սերունդ են թողել
      \end{card}}
      \onslide<4->{%
      \begin{card}[colback=butter!45!cream]
        \emoji[1.3em]{brain}\ {\blackfont 4. Ուղեղի չափ}\par\smallskip
        Բարձրահասակների ուղեղն ավելի մեծ է, իսկ ուղեղի չափն ու խելքը թույլ կապ ունեն
      \end{card}}
    \end{column}
  \end{columns}
\end{frame}

% Sexual selection and its problems
\begin{frame}{Սեռական ընտրության վարկածի խնդիրները}
  \begin{columns}[T]
    \begin{column}{0.56\textwidth}
      \begin{itemize}
        \item Վարկածը շատ \textbf{տղամարդակենտրոն} է. իբր մի երկու հատկանիշ ունեցիր, ու կանայք ինքնաբերաբար կգան
        \item<2-> Բարձրահասակ տղամարդիկ ունենում են նաև բարձրահասակ \textbf{աղջիկներ}, իսկ շատ տղամարդիկ վախենում են բարձրահասակ ու խելացի կանանցից
        \item<3-> Վստահությունը գրավիչ է, իսկ խելացիները հաճախ ավելի քիչ վստահ են։ Ու «nerd», «geek» բառերը երկար ժամանակ վիրավորանք էին
      \end{itemize}
    \end{column}
    \begin{column}{0.38\textwidth}
      \centering
      \bigemoji[1.2cm]{giraffe}\ \bigemoji[1.2cm]{nerd}\par\medskip
      \begin{bookquote}
        ...ինչպես ցեցերը դեպի նիհարավուն, սրամիտ կրակը։
      \end{bookquote}
      {\scriptsize Բըրնեթ. «Այդպես են ասում իմ բարձրահասակ ու խելացի ընկերուհիները։ Իսկ իմ բոլոր ընկերուհիները խելացի են»։\par}
    \end{column}
  \end{columns}
\end{frame}

% Nutrition, health and brain size
\begin{frame}{Սնունդ, առողջություն, ուղեղի չափ}
  \begin{columns}[T]
    \begin{column}{0.5\textwidth}
      \begin{card}[colback=mint!35!cardwhite]
        \emoji[1.3em]{apple}\ Լավ սնունդն ու առողջությունը կարող են մեծացնել և՛ հասակը, և՛ խելքը։ Բայց շատերը բարեկեցիկ կյանքով էլ մնում են ցածրահասակ։ Կամ ապուշ։ Կամ երկուսը։
      \end{card}
    \end{column}
    \begin{column}{0.46\textwidth}
      \onslide<2->{%
      \begin{card}[colback=butter!45!cream]
        \emoji[1.3em]{brain}\ Ավելի խելացիների մոտ նախաճակատային կեղևն ու հիպոկամպը ավելի մեծ են։ Մեծ ուղեղը \textbf{հնարավորություն} է, ոչ թե երաշխիք։
      \end{card}}
    \end{column}
  \end{columns}
  \onslide<3->{%
  \begin{bookquote}
    Թանկ սպորտային կոշիկը քեզ ավելի արագ չի դարձնում, բայց կարող է դրդել, որ դառնաս։ Նույնը կարելի է ասել գեների մասին։
  \end{bookquote}}
\end{frame}

% Intelligence and culture: fish and water
\begin{frame}{Ձուկն ու ջուրը}
  \begin{columns}[c]
    \begin{column}{0.5\textwidth}
      Խելքի վրա ազդում են...\par\medskip
      \chip{dna}{գեներ}\ \chip{family}{դաստիարակություն}\par\smallskip
      \chip{graduation}{կրթության որակ}\ \chip{earth}{մշակույթ}\par\smallskip
      \chip{speech}{կարծրատիպեր}\ \chip{stethoscope}{առողջություն}\par\smallskip
      \chip{palette}{հետաքրքրություններ}\ \chip{hospital}{խանգարումներ}
    \end{column}
    \begin{column}{0.46\textwidth}
      \centering
      \bigemoji[1.3cm]{fish}\ \bigemoji[1.3cm]{wave}\par\medskip
      \onslide<2->{%
      \begin{bookquote}
        Մարդու խելքը մշակույթից չես առանձնացնի, ինչպես ձկան զարգացումը՝ ջրից։ Իսկ եթե առանձնացնես, ձկան զարգացումը շատ «կարճ» կլինի։
      \end{bookquote}}
    \end{column}
  \end{columns}
\end{frame}

% The Kpelle and "the way a fool would do it"
\begin{frame}{Կպելլեները և «հիմարի պես»}
  \begin{columns}[c]
    \begin{column}{0.58\textwidth}
      \begin{enumerate}
        \item<1-> Մայքլ Քոուլի խումբը գնաց Լիբերիա՝ \textbf{կպելլե} ժողովրդի մոտ, որ արևմտյան կրթության հետ գրեթե շփում չուներ
        \item<2-> Խնդիր. առարկաները խմբավորել։ Կպելլեները «ձախողում» էին, նույնիսկ հուշումներով
        \item<3-> Հետազոտողները սկսեցին մտածել, թե նրանք պակաս խելացի են
        \item<4-> Հուսահատված հետազոտողը խնդրեց անել այնպես, \textbf{ինչպես հիմարը կաներ}։ Ու նրանք անմիջապես տվեցին «ճիշտ» պատասխանները
      \end{enumerate}
    \end{column}
    \begin{column}{0.36\textwidth}
      \centering
      \bigemoji[1.3cm]{earth}\par\medskip
      \onslide<4->{\mascot[2.2cm]{smug}}
    \end{column}
  \end{columns}
\end{frame}

% Two ways to sort
\begin{frame}{Ո՞վ է իրականում «հիմարը»}
  \begin{columns}[T]
    \begin{column}{0.47\textwidth}
      \begin{card}[colback=lavender!35!cardwhite]
        {\blackfont Հետազոտողների «խելացի» ձևը}\par
        ըստ \textbf{դասակարգի}\par\medskip
        \emoji{axe}\emoji{knife}\ գործիքներ\par\smallskip
        \emoji{sheep}\emoji{fish}\ կենդանիներ\par\smallskip
        \emoji{wood}\ փայտից իրեր
      \end{card}
    \end{column}
    \begin{column}{0.47\textwidth}
      \onslide<2->{%
      \begin{card}[colback=mint!35!cardwhite]
        {\blackfont Կպելլեների ձևը}\par
        ըստ \textbf{գործառույթի}\par\medskip
        \emoji{potato}\emoji{fish}\ ինչ կարող եմ ուտել\par\smallskip
        \emoji{jeans}\ ինչ կարող եմ հագնել\par\smallskip
        \emoji{axe}\ ինչով կարող եմ փորել, կտրել
      \end{card}}
    \end{column}
  \end{columns}
  \vspace{2mm}
  \onslide<3->{%
  \begin{card}[colback=butter!45!cream]
    Հողից ապրողի համար վերացական դասակարգերը անիմաստ են. դա \textbf{հիմարի} գործ է։ Խելքը, որ չափում ենք, կախված է մշակույթից ու մեր նախապաշարմունքներից։
  \end{card}}
\end{frame}

% The Pygmalion effect (+ science update)
\begin{frame}{Պիգմալիոնի էֆեկտը}
  \begin{columns}[T]
    \begin{column}{0.54\textwidth}
      \begin{itemize}
        \item 1965, Ռոզենթալ և Ջեյկոբսոն. ուսուցիչներին ասացին, որ որոշ աշակերտներ \textbf{շնորհալի} են։ Իրականում սովորական էին
        \item<2-> Ու այդ աշակերտներն իսկապես ավելի լավ արդյունքներ ցույց տվեցին
        \item<3-> Ուսանողները, որոնց ասում էին՝ խելքը \textbf{ֆիքսված} է, ավելի վատ էին գրում, իսկ «փոփոխական է» լսածները՝ ավելի լավ
      \end{itemize}
    \end{column}
    \begin{column}{0.42\textwidth}
      \onslide<4->{%
      \begin{scienceupdate}
        Էֆեկտներն իրական են, բայց \textbf{ավելի փոքր}, քան սկզբում էր թվում։ «Խելքը զարգանում է» մտածելակերպի ծրագրերը միջինում շատ փոքր ազդեցություն ունեն, մի փոքր ավելի մեծ՝ դժվարացող ուսանողների մոտ։
      \end{scienceupdate}}
    \end{column}
  \end{columns}
  \source{Jussim \& Harber, Pers. Soc. Psychol. Rev., 2005; Sisk et al., Psychol. Sci., 2018}
\end{frame}

% Burnett apologises
\begin{frame}{Ներողություն, իմ մեղքն է}
  \begin{columns}[c]
    \begin{column}{0.56\textwidth}
      Գուցե դրա համար էլ են բարձրահասակներն ավելի խելացի. եթե շուտ ես մեծանում, քեզ \textbf{ավելի մեծ} են համարում, հետդ ավելի լուրջ են խոսում, ու զարգացող ուղեղդ «հասնում է» այդ սպասումներին։\par\medskip
      \onslide<2->{Ամեն դեպքում, \textbf{ինքդ քեզ հավատալը} կարևոր է։\par}
    \end{column}
    \begin{column}{0.4\textwidth}
      \centering
      \onslide<3->{%
      \mascot[1.5cm]{sad}\par\smallskip
      \begin{bookquote}
        Ամեն անգամ, երբ գրել եմ, թե խելքը «ֆիքսված» է, խանգարել եմ ձեր զարգացմանը։ Ներողություն, իմ մեղքն է։
      \end{bookquote}}
    \end{column}
  \end{columns}
\end{frame}

% The Flynn effect (+ science update)
\begin{frame}{Ֆլինի էֆեկտը}
  \begin{columns}[T]
    \begin{column}{0.54\textwidth}
      \begin{itemize}
        \item IQ թեստերի արդյունքները աշխարհով մեկ \textbf{աճում են} ամեն սերնդի հետ՝ և՛ հեղուկ, և՛ բյուրեղացած
        \item<2-> Ինչու՞։ Գուցե ավելի լավ կրթություն, բժշկություն, տեղեկության ու տեխնոլոգիաների հասանելիություն...
        \item<3-> ...կամ քնած մուտանտային ունակություններ, որոնք մեզ հանճարների հասարակություն կդարձնեն։ Դրա ապացույց չկա, բայց լավ ֆիլմ կստացվեր \emoji{film}
      \end{itemize}
    \end{column}
    \begin{column}{0.42\textwidth}
      \onslide<4->{%
      \begin{scienceupdate}
        1990-ականներից Նորվեգիայում, Դանիայում ու Ֆինլանդիայում աճը կանգնեց ու նույնիսկ \textbf{նվազեց}։ Նորվեգացի եղբայրների տվյալները ցույց տվեցին, որ պատճառը \textbf{միջավայրն} է, ոչ թե գեները։
      \end{scienceupdate}}
    \end{column}
  \end{columns}
  \source{Bratsberg \& Rogeberg, PNAS, 2018}
\end{frame}

% The chapter's conclusion
\begin{frame}{Եվ ի վերջո...}
  \begin{columns}[c]
    \begin{column}{0.52\textwidth}
      Խելքը դժվար է սահմանել ու չափել, բայց այն \textbf{կա}, ու կարելի է ուսումնասիրել։\par\medskip
      \onslide<2->{Մեկ ընդհանուր ունակություն է՝ կազմված մի քանիսից։ Ու կարևորը, հավանաբար, ուղեղի մասերի \textbf{կապերն} են։\par\medskip}
      \onslide<3->{Գեներն ու դաստիարակությունը շատ բան են որոշում... բայց եթե պատրաստ ես աշխատել, խելքը կարելի է բարձրացնել։ \textbf{Գուցե}։\par}
    \end{column}
    \begin{column}{0.44\textwidth}
      \centering
      \onslide<4->{%
      \mascot[1.4cm]{happy}\par\smallskip
      \begin{bookquote}
        Խելքն ուսումնասիրելը նման է առանց ձևանմուշի սվիտեր գործելուն՝ բրդի փոխարեն շաքարավատով։ Զարմանալի է, որ փորձում ենք։
      \end{bookquote}}
    \end{column}
  \end{columns}
\end{frame}
